% 分野別類題: 微分計算（逆三角関数） 解答例
% コンパイル: TEXINPUTS="../../../00_common:$TEXINPUTS" latexmk -xelatex answers.tex
% 逆三角関数の微分公式（arcsin, arccos, arctan）の導出と、合成・積の微分計算5問の解答。
\documentclass[a4paper,11pt]{article}
\usepackage{preamble}
\usepackage{shoakubox}

\begin{document}

\kamokumei{類題演習 解答例 --- 微分計算（逆三角関数）}

\daimon{【解法】逆三角関数の微分公式とその導出}

\noindent
逆関数の微分法によれば、$y = f(x)$ が単調で微分可能、$f'(x) \neq 0$ のとき、逆関数 $x = f^{-1}(y)$ も微分可能であり
\[
\frac{dx}{dy} = \frac{1}{\dfrac{dy}{dx}}
\]
が成り立つ。この原理を用いて、逆三角関数 $\arcsin x,\ \arccos x,\ \arctan x$ の微分公式を導く。

\medskip
\noindent
\textbf{(1) $\arcsin x$ の導出.}\quad
$y = \arcsin x$（$-1 < x < 1$、$-\dfrac{\pi}{2} < y < \dfrac{\pi}{2}$）とおくと、$x = \sin y$ である。両辺を $x$ について微分すると
\[
1 = \cos y \cdot \frac{dy}{dx}
\quad\Longrightarrow\quad
\frac{dy}{dx} = \frac{1}{\cos y}
\]
$-\dfrac{\pi}{2} < y < \dfrac{\pi}{2}$ では $\cos y > 0$ だから、$\cos y = \sqrt{1-\sin^{2} y} = \sqrt{1-x^{2}}$ となり
\[
\frac{d}{dx}\arcsin x = \frac{1}{\sqrt{1-x^{2}}}
\]

\medskip
\noindent
\textbf{(2) $\arccos x$ の導出.}\quad
$y = \arccos x$（$0 < y < \pi$）とおくと $x = \cos y$。両辺を $x$ で微分すると
\[
1 = -\sin y \cdot \frac{dy}{dx}
\quad\Longrightarrow\quad
\frac{dy}{dx} = -\frac{1}{\sin y}
\]
$0 < y < \pi$ では $\sin y > 0$ だから $\sin y = \sqrt{1-\cos^{2} y} = \sqrt{1-x^{2}}$ となり
\[
\frac{d}{dx}\arccos x = -\frac{1}{\sqrt{1-x^{2}}}
\]
（別解として $\arccos x = \dfrac{\pi}{2} - \arcsin x$ を微分しても同じ結果が得られる。）

\medskip
\noindent
\textbf{(3) $\arctan x$ の導出.}\quad
$y = \arctan x$（$-\dfrac{\pi}{2} < y < \dfrac{\pi}{2}$）とおくと $x = \tan y$。両辺を $x$ で微分すると
\[
1 = \sec^{2} y \cdot \frac{dy}{dx} = (1+\tan^{2} y)\frac{dy}{dx}
\quad\Longrightarrow\quad
\frac{dy}{dx} = \frac{1}{1+\tan^{2} y} = \frac{1}{1+x^{2}}
\]
よって
\[
\frac{d}{dx}\arctan x = \frac{1}{1+x^{2}}
\]

\medskip
\noindent
これらの公式に合成関数の微分法（連鎖律）を組み合わせることで、$\arcsin u(x),\ \arccos u(x),\ \arctan u(x)$ の微分は
\[
\frac{d}{dx}\arcsin u = \frac{u'}{\sqrt{1-u^{2}}}, \qquad
\frac{d}{dx}\arccos u = -\frac{u'}{\sqrt{1-u^{2}}}, \qquad
\frac{d}{dx}\arctan u = \frac{u'}{1+u^{2}}
\]
となる。以下の問2以降ではこれらを利用する。

\clearpage
\begin{mondaiboxS}{問1}
$y = \arcsin x$（$-1 < x < 1$）について、逆関数の微分法を用いて $\dfrac{dy}{dx}$ を導出せよ。
\end{mondaiboxS}
\kaitou
$x = \sin y$（$-\dfrac{\pi}{2} < y < \dfrac{\pi}{2}$）とおいて両辺を $x$ で微分すると
\[
1 = \cos y \cdot \frac{dy}{dx}
\]
この範囲で $\cos y > 0$ なので $\cos y = \sqrt{1-\sin^{2}y} = \sqrt{1-x^{2}}$。したがって
\[
\boxed{\; \frac{dy}{dx} = \frac{1}{\sqrt{1-x^{2}}} \;}
\]
（検算: $x=0$ のとき $y=0$ で $\dfrac{dy}{dx} = \dfrac{1}{\sqrt{1-0}} = 1$。実際 $y=\arcsin x$ は原点付近で $y \approx x$ なので傾き $1$ と整合する。）

\clearpage
\begin{mondaiboxS}{問2}
$y = \arctan x$ について、逆関数の微分法を用いて $\dfrac{dy}{dx}$ を導出せよ。
\end{mondaiboxS}
\kaitou
$x = \tan y$（$-\dfrac{\pi}{2} < y < \dfrac{\pi}{2}$）とおいて両辺を $x$ で微分すると
\[
1 = \sec^{2} y \cdot \frac{dy}{dx} = (1+\tan^{2}y)\frac{dy}{dx} = (1+x^{2})\frac{dy}{dx}
\]
したがって
\[
\boxed{\; \frac{dy}{dx} = \frac{1}{1+x^{2}} \;}
\]
（検算: $x=1$ のとき $y = \arctan 1 = \dfrac{\pi}{4}$ で傾きは $\dfrac{1}{1+1}=\dfrac{1}{2}$。$\tan y$ の逆関数として、$x=1$ 付近で $\tan y$ の傾き $\sec^{2}(\pi/4)=2$ の逆数 $\dfrac12$ と一致する。）

\clearpage
\begin{mondaiboxS}{問3}
次の関数を微分せよ。 $y = \arccos(3x)$
\end{mondaiboxS}
\kaitou
$u = 3x$ とおくと $y = \arccos u$、$u' = 3$。合成関数の微分公式 $\dfrac{d}{dx}\arccos u = -\dfrac{u'}{\sqrt{1-u^{2}}}$ より
\[
\frac{dy}{dx} = -\frac{3}{\sqrt{1-(3x)^{2}}}
\]
したがって
\[
\boxed{\; \frac{dy}{dx} = -\frac{3}{\sqrt{1-9x^{2}}} \quad \left(-\frac{1}{3} < x < \frac{1}{3}\right) \;}
\]
（検算: $x=0$ のとき $y=\arccos 0 = \pi/2$。$h$ を微小量として $y(h) = \arccos(3h) \approx \dfrac{\pi}{2} - 3h$（$\because \arccos(\varepsilon)\approx \pi/2-\varepsilon$）であり、傾き $-3$ は上式に $x=0$ を代入した値 $-3/\sqrt{1-0}=-3$ と一致する。）

\clearpage
\begin{mondaiboxS}{問4}
次の関数を微分せよ。 $y = \arctan\!\left(\dfrac{2x}{1-x^{2}}\right)$
\end{mondaiboxS}
\kaitou
$u = \dfrac{2x}{1-x^{2}}$ とおく。商の微分法より
\[
u' = \frac{2(1-x^{2}) - 2x\cdot(-2x)}{(1-x^{2})^{2}}
= \frac{2 - 2x^{2} + 4x^{2}}{(1-x^{2})^{2}}
= \frac{2 + 2x^{2}}{(1-x^{2})^{2}}
= \frac{2(1+x^{2})}{(1-x^{2})^{2}}
\]
また
\[
1 + u^{2} = 1 + \frac{4x^{2}}{(1-x^{2})^{2}}
= \frac{(1-x^{2})^{2} + 4x^{2}}{(1-x^{2})^{2}}
= \frac{1 - 2x^{2} + x^{4} + 4x^{2}}{(1-x^{2})^{2}}
= \frac{(1+x^{2})^{2}}{(1-x^{2})^{2}}
\]
したがって、公式 $\dfrac{d}{dx}\arctan u = \dfrac{u'}{1+u^{2}}$ より
\[
\frac{dy}{dx} = \frac{\dfrac{2(1+x^{2})}{(1-x^{2})^{2}}}{\dfrac{(1+x^{2})^{2}}{(1-x^{2})^{2}}}
= \frac{2(1+x^{2})}{(1+x^{2})^{2}}
= \frac{2}{1+x^{2}}
\]
（$x^{2} \neq 1$、すなわち $x \neq \pm 1$ で定義される。）したがって
\[
\boxed{\; \frac{dy}{dx} = \frac{2}{1+x^{2}} \;}
\]
（検算: これは $\arctan(2x/(1-x^{2})) = 2\arctan x$（$|x|<1$、正接の2倍角公式 $\tan(2\theta) = \dfrac{2\tan\theta}{1-\tan^{2}\theta}$ に $\theta=\arctan x$ を代入した恒等式）と一致することからも確認できる。実際 $y=2\arctan x$ を直接微分すると $\dfrac{dy}{dx} = 2\cdot\dfrac{1}{1+x^{2}} = \dfrac{2}{1+x^{2}}$ となり、上の結果と一致する。）

\clearpage
\begin{mondaiboxS}{問5}
次の関数を微分せよ。またその結果をできるだけ簡単な形にせよ。 $y = x\arcsin x + \sqrt{1-x^{2}}$
\end{mondaiboxS}
\kaitou
積の微分法と合成関数の微分法を用いる。第1項 $x\arcsin x$ は
\[
\frac{d}{dx}\bigl(x\arcsin x\bigr) = 1\cdot \arcsin x + x\cdot \frac{1}{\sqrt{1-x^{2}}}
= \arcsin x + \frac{x}{\sqrt{1-x^{2}}}
\]
第2項 $\sqrt{1-x^{2}} = (1-x^{2})^{1/2}$ は
\[
\frac{d}{dx}\sqrt{1-x^{2}} = \frac{1}{2}(1-x^{2})^{-1/2}\cdot(-2x) = -\frac{x}{\sqrt{1-x^{2}}}
\]
したがって
\[
\frac{dy}{dx} = \left(\arcsin x + \frac{x}{\sqrt{1-x^{2}}}\right) + \left(-\frac{x}{\sqrt{1-x^{2}}}\right) = \arcsin x
\]
すなわち
\[
\boxed{\; \frac{dy}{dx} = \arcsin x \;}
\]
（検算: $x=0$ のとき $y(0) = 0\cdot 0 + 1 = 1$、$y'(0) = \arcsin 0 = 0$。実際 $y=x\arcsin x+\sqrt{1-x^{2}}$ は $x=0$ で極値（極小値 $1$）をとる偶関数であり、$y'(0)=0$ と整合する。さらに一般の $x$ で $\dfrac{x}{\sqrt{1-x^2}}$ の項が打ち消し合う構造になっており、計算過程にも矛盾はない。）

\end{document}
